
  \documentclass{lxaiproposal}
  
  
%   \usepackage[english,french]{babel}   % "babel.sty" + "french.sty"
   \usepackage[english]{babel}   % "babel.sty" + "french.sty"
   
% \usepackage[english,francais]{babel} % "babel.sty"
% \usepackage{french}                  % "french.sty"
  \usepackage{times}			% ajout times le 30 mai 2003
 
\usepackage{epsfig}
\usepackage{graphicx}
\usepackage{amsmath}
\usepackage{amssymb}
 \usepackage{booktabs}
 \usepackage{multirow}
\usepackage{pifont}
\usepackage{caption}
\usepackage{subcaption}



\usepackage{array}
\usepackage{color}
\usepackage{colortbl}

\usepackage{pifont}
\usepackage{amssymb}
\usepackage{latexsym}

\usepackage{booktabs}

%% --------------------------------------------------------------
%% FONTS CODING ?
% \usepackage[OT1]{fontenc} % Old fonts
% \usepackage[T1]{fontenc}  % New fonts (preferred)
%% ==============================================================

\title{TITLE PROPOSAL}

\author{\coord{PI FirstNAME}{PI FamilyNAME}{1} \\
       \coord{CO-PI FirstNAME}{CO-PI FamilyNAME}{2}
      }

\address{\affil{1}{Affiliation 1, Address}
         \affil{2}{Affiliation 2, Address}}

%% If all authors have the same address %%%%%%%%%%%%%%%%%%%%%%%%%%%%%%%%%%%%%%%
%                                                                             %
%   \auteur{\coord{Michel}{Dupont}{},                                         %
%           \coord{Marcel}{Dupond}{},                                         %
%           \coord{Michelle}{Durand}{},                                       %
%           \coord{Marcelle}{Durand}{}}                                       %
%                                                                             %
%   \adress{\affil{}{Laboratoire Traitement des Signaux et des Images \\      %
%     1 rue de la Science, BP 00000, 99999 Nouvelleville Cedex 00, France}}   %
%                                                                             %
%                                                                             %
%%%%%%%%%%%%%%%%%%%%%%%%%%%%%%%%%%%%%%%%%%%%%%%%%%%%%%%%%%%%%%%%%%%%%%%%%%%%%%%

\email{PI email, CO-PI email}

\englishabstract{}

\begin{document}
\maketitle
%
\section{Background and Motivation}
\vspace*{-3mm}

Provide a problem statement and discuss the motivation for carrying out the project. 
\textbf{(1-2 paragraphs)}

\section{Relevance and related works}
\vspace*{-3mm}

You should describe Intellectual Merit: The intellectual Merit criterion encompasses the potential to advance knowledge (Technical relevance for startups/NGOs, Scientific relevance for academic projects) as compared to related works.

Discuss in a more focused manner how your AI project will solve the problem stated above and how you plan to build upon the state of the art.
\textbf{(1-2 paragraphs)}


\section{Methods and Approach}
\vspace*{-3mm}

Describe how you plan to solve the problem, breaking the proposed solution into specific items, i.e., stages with well defined goals, milestones, measures of success, deliverables, etc. 

Please include estimated milestones with timeline estimates, such as for datasets, code releases, technical reports, publications, applications, presentations, etc. 
\textbf{(1 page)}


\section{Expected Results and Impact}
\vspace*{-3mm}

Describe which impact and outcomes your project will have. For academic projects you can discuss in terms of human resources formation, development of any publicly available dataset, new capabilities, impact in larger projects, advance state-of-the-art knowledge, etc. 
\textbf{(1 page)}

\subsection{Broader Impacts}


The Broader Impacts criterion encompasses the potential to benefit society and contribute to the achievement of specific, desired societal outcomes. Consider the following questions.


1. What is the potential for the proposed activity to benefit society or advance desired societal outcomes?

2. To what extent do the proposed activities suggest and explore creative, original or potentially transformative concepts?

3. Is the plan for carrying out the proposed activities well-reasoned, well-organized and based on sound rationale? Does the plan incorporate a mechanism to assess success?

4. How well qualified is the individual, team or institution to conduct the proposed activities?

5. Are there adequate resources available to the principal investigator (either at the home institution or through collaborations) to carry out the proposed activities

6. How will the access to the computer resources provided by the consortium help you advance your research if your project is supported?

\textbf{(1-2 paragraphs)}


\section{Data policy}
\vspace*{-3mm}


Our goal is to support work where the output will be made available to the broader research community. To that end, we ask that you provide us with a few sentences sharing what you intend to do with the output of your project (e.g. open sourcing code, making data sets public, etc.).

Please note that there are no legal requirements once a project is selected This is simply a statement of your current intentions. 
\textbf{(1 paragraph)} \\

Examples of cite, table and figure: This is a cite \cite{daudon2018recurrence}. Table \ref{tab:my-table}. Figure \ref{fig:my-fig}.


% Please add the following required packages to your document preamble:
\begin{table}[h!]
\centering
\caption{A typical table.}
\label{tab:my-table}
\begin{tabular}{@{}ccccc@{}}
\toprule
\textbf{Items/Cols} & \textbf{Col 1} & \textbf{Col 2} & \textbf{Col 3} & \textbf{Col 4} \\ \midrule
\textbf{Item 1}     & Value          & Value          & Value          & Value          \\
\textbf{Item 2}     & Value          & Value          & Value          & Value          \\
\textbf{Item 3}     & Value          & Value          & Value          & Value          \\ \bottomrule
\end{tabular}
\end{table}

 %\begin{figure} [h!]
 %    \centering
 %    \includegraphics[width=1\linewidth]{images/fig1.png}
 %  \caption{A typical figure.}
 %    \label{fig:my-fig}
 %    \end{figure}

\bibliographystyle{ieee_fullname}
\bibliography{references}

%\begin{thebibliography}{99}
%\end{thebibliography}

\end{document}
